\documentclass[11pt]{amsart}

\usepackage[T1]{fontenc}
\usepackage{lmodern}
\usepackage{microtype}
\usepackage{amsmath,amssymb,amsthm}
\usepackage[hidelinks]{hyperref}

\hypersetup{
  pdftitle={A Non-Cayley Vertex-Transitive 3-Rainbow Domination Regular Graph},
  pdfsubject={Rainbow domination in cubic vertex-transitive graphs},
  pdfkeywords={rainbow domination, Cayley graph, vertex-transitive graph, cubic graph}
}

\newtheorem{theorem}{Theorem}
\theoremstyle{remark}
\newtheorem{remark}[theorem]{Remark}

\newcommand{\Aut}{\operatorname{Aut}}
\newcommand{\gr}{\gamma_{r3}}

\title[A non-Cayley vertex-transitive 3-RDR graph]
{A Non-Cayley Vertex-Transitive\\
3-Rainbow Domination Regular Graph}

\date{}

\subjclass[2020]{05C69, 05C25}
\keywords{rainbow domination, Cayley graph, vertex-transitive graph, cubic graph}

\begin{document}

\begin{abstract}
We exhibit a connected cubic vertex-transitive graph $X$ of order $108$ such
that $\gr(X)=54$ and $X$ is not a Cayley graph.  The graph is
$\mathrm{CVT}[108,42]$ in the cubic vertex-transitive census of
Poto\v{c}nik, Spiga, and Verret.  Its optimal $3$-rainbow dominating function
is explicit, and its non-Cayley property follows from an exact computation of
its full automorphism group.  This answers the non-Cayley existence question
in Kuzman's Question~\textup{(Q5)} for $d=3$.
\end{abstract}

\maketitle

A \emph{$3$-rainbow dominating function} on a graph $G$ is a map
$f\colon V(G)\to 2^{\{1,2,3\}}$ such that
\[
 f(v)=\varnothing
 \quad\Longrightarrow\quad
 \bigcup_{u\in N(v)}f(u)=\{1,2,3\}.
\]
Its weight is $w(f)=\sum_{v\in V(G)}|f(v)|$, and the minimum possible
weight is the $3$-rainbow domination number $\gr(G)$.  A cubic graph $G$ is
\emph{$3$-rainbow domination regular}, or \emph{$3$-RDR}, when
$\gr(G)=|V(G)|/2$.  Kuzman asked whether a non-Cayley vertex-transitive
$3$-RDR graph exists~\cite[Question~\textup{(Q5)}]{Kuzman2025}.

\begin{theorem}\label{thm:main}
Let $X$ be the graph specified in Appendix~\ref{app:witness}.  Then $X$ is a
connected cubic vertex-transitive non-Cayley graph of order $108$, and
\[
 \gr(X)=54.
\]
Consequently, non-Cayley vertex-transitive $3$-RDR graphs exist.
\end{theorem}

\begin{proof}
Decoding the sparse6 string in Appendix~\ref{app:witness} gives a connected
cubic graph on $108$ vertices.  It is isomorphic to the census graph
$\mathrm{CVT}[108,42]$~\cite{PSV2013,CensusData}.

Set
\begin{align*}
C_1={}&\{4,5,18,19,23,33,52,54,60,61,67,74,78,79,83,86,104,107\},\\
C_2={}&\{8,10,11,12,21,35,40,55,58,63,65,72,80,82,84,98,99,101\},\\
C_3={}&\{1,20,24,26,31,34,37,38,39,41,50,56,76,85,88,97,100,103\}.
\end{align*}
These sets are pairwise disjoint and each has size $18$.  Every vertex outside
$C_1\cup C_2\cup C_3$ has exactly one neighbor in each $C_i$.  Hence the
assignment
\[
 f(v)=
 \begin{cases}
  \{i\},&v\in C_i,\\
  \varnothing,&v\notin C_1\cup C_2\cup C_3
 \end{cases}
\]
is a $3$-rainbow dominating function of weight $54$.

For completeness, let $g$ be any $3$-rainbow dominating function on a cubic
graph $G$ of order $n$, let $w=w(g)$, and put
$V_0=\{v:g(v)=\varnothing\}$.  Each vertex of $V_0$ receives all three colors
from its neighborhood, while each unit of weight is counted at most three
times.  Thus
\[
 3|V_0|
 \leq \sum_{v\in V_0}\sum_{u\in N(v)}|g(u)|
 \leq 3w.
\]
Also $|V(G)\setminus V_0|\leq w$, so $n\leq 2w$.  Therefore
$\gr(G)\geq n/2$, and the displayed function proves $\gr(X)=54$.

It remains to verify the symmetry assertions.  An exhaustive enumeration of
all adjacency-preserving permutations gives, for $A=\Aut(X)$,
\[
 |A|=216,
 \qquad A\text{ is transitive},
 \qquad |[A,A]|=54,
 \qquad A/[A,A]\cong C_2\times C_2.
\]
Consequently, $A$ has exactly three subgroups of index $2$.  Each has two
vertex orbits, both of size $54$.  If $X$ were a Cayley graph, $A$ would
contain a regular subgroup of order $108$; such a subgroup would have index
$2$ in $A$ and would be transitive, a contradiction.  Hence $X$ is
non-Cayley.

\end{proof}

\begin{remark}
Theorem~\ref{thm:main} settles the non-Cayley existence question in
Kuzman's Question~\textup{(Q5)} for $d=3$.  No claim is made here about
minimum order or the cases $d\geq4$.
\end{remark}

\appendix
\section{The witness}\label{app:witness}

Concatenate the following four strings without whitespace and decode the
result in standard sparse6 format.  Number the decoded vertices
$0,1,\ldots,107$ and then add $1$ to every label.

\begin{center}
\footnotesize
\verb+:~?@k_GAA_WAC`WM@_gWGaWiF_wQDb?yNcHEL`g]GaWgPcpOUd`]Wc`e+\\
\verb+YexUJbGuMGw}OEhw_GiGchY]ggXMSdX]WeQ]jJAymJzApkjMsjYE\fh}+\\
\verb+_grIcOIS}hbMvMRqwMbu{Njs~oKFAoz`CpZlNjYijjI{okTHRqcpOqst+\\
\verb+PsTJLTKlTmzXHlzeyVz|@VlhZVLp\rK|bqDfIre@cstPfvd|dxEXf+
\end{center}

\begin{thebibliography}{9}

\bibitem{Kuzman2025}
B.~Kuzman,
\emph{On cubic rainbow domination regular graphs},
Discrete Appl. Math. \textbf{373} (2025), 26--38.
\href{https://doi.org/10.1016/j.dam.2025.04.046}
{doi:10.1016/j.dam.2025.04.046}.

\bibitem{PSV2013}
P.~Poto\v{c}nik, P.~Spiga, and G.~Verret,
\emph{Cubic vertex-transitive graphs on up to 1280 vertices},
J. Symbolic Comput. \textbf{50} (2013), 465--477.
\href{https://doi.org/10.1016/j.jsc.2012.09.002}
{doi:10.1016/j.jsc.2012.09.002}.

\bibitem{CensusData}
J.~Azarija, K.~Ber\v{c}i\v{c}, F.~Koprivec, T.~Marc, P.~Poto\v{c}nik,
P.~Rotar, P.~Spiga, G.~Verret, and J.~Vidali,
\emph{Cubic vertex-transitive graphs on up to 1280 vertices},
Zenodo, version 1 (2022).
\href{https://doi.org/10.5281/zenodo.6576526}
{doi:10.5281/zenodo.6576526}.

\end{thebibliography}

\end{document}